% Part: proof-theory
% Chapter: normalization
% Section: segments

\documentclass[../../../include/open-logic-section]{subfiles}

\begin{document}

\olfileid{pt}{nor}{seg}

\olsection{قطعه‌ها و برش‌ها}

اگر قاعده‌ای !!^a{introduction} و سپس قاعده‌ای !!a{elimination} به کار
رود، در !!a{proof} یک مسیر انحرافی پدید می‌آید؛ بهنجار کردن !!a{proof}
یعنی حذف چنین مسیرهایی. بااین‌حال، قواعد \Elim{\lor} و
\Elim{\lexists} حتی تعریف «مسیر انحرافی‌ای که می‌خواهیم حذف کنیم» را
کمی پیچیده می‌کنند. مسئله این است که، به‌سبب این قواعد، تضمینی نیست
که در یک مسیر انحرافی قاعدهٔ !!{elimination} \emph{بی‌درنگ} پس از
قاعدهٔ !!{introduction} بیاید. ممکن است میان آن‌ها یک \Elim{\lor} یا
\Elim{\lexists} قرار گیرد؛ برای نمونه:
\[
\AxiomC{}
\DeduceC{$\lexists[x][!C(x)]$}
\AxiomC{$\Discharge{!A}{x}$}
\DeduceC{$!B$}
\DischargeRule{\Intro{\lif}}{x}
\UnaryInfC{$!A \lif !B$}
\DischargeRule{\Elim{\lexists}}{y}
\BinaryInfC{$!A \lif !B$}
\AxiomC{}
\DeduceC{$!A$}
\RightLabel{\Elim{\lif}}
\BinaryInfC{$!B$}
\DisplayProof
\]
در واقع، هر تعداد قاعدهٔ \Elim{\lor} یا \Elim{\lexists} می‌تواند
\Intro{\lif} را از \Elim{\lif} جدا کند؛ برای نمونه:
\[
\AxiomC{}
\DeduceC{$\lexists[x][!C(x)]$}
\AxiomC{}
\DeduceC{$!D \lor !E$}
\AxiomC{$\Discharge{!A}{x}$}
\DeduceC{$!B$}
\DischargeRule{\Intro{\lif}}{x}
\UnaryInfC{$!A \lif !B$}
\AxiomC{$\Discharge{!A}{x}$}
\DeduceC{$!B$}
\DischargeRule{\Intro{\lif}}{x}
\UnaryInfC{$!A \lif !B$}
\RightLabel{\Elim{\lor}}
\TrinaryInfC{$!A \lif !B$}
\DischargeRule{\Elim{\lexists}}{y}
\BinaryInfC{$!A \lif !B$}
\AxiomC{}
\DeduceC{$!A$}
\RightLabel{\Elim{\lif}}
\BinaryInfC{$!B$}
\DisplayProof
\]
این مثال نشان می‌دهد که ممکن است یک \Elim{\lif} واحد در چند مسیر
انحرافی دخیل باشد، زیرا پس از بیش از یک~\Intro{\lif} می‌آید. برای
درنظرگرفتن همهٔ این امکان‌ها، مسیرهای انحرافی را با استفاده از گونه‌ای
دنبالهٔ وقوع‌های !!{formula} در یک برهان \Log{N1i} تعریف می‌کنیم
که به آن‌ها \emph{قطعه} می‌گوییم. در مثال ما، این دنباله‌ای از وقوع‌های
~$!A \lif !B$ است که در آن نتیجهٔ هر $\Elim{\lor}$ یا
$\Elim{\lexists}$ در ادامهٔ مقدمهٔ فرعی آن می‌آید. در این مثال دو چنین
قطعه‌ای وجود دارد و هریک سه وقوع از~$!A \lif !B$ را در بر دارد. تعریف
رسمی چنین است.

\begin{defn}
یک \emph{قطعه} در !!{proof} \Log{N1i}ای~$\delta$ دنباله‌ای از وقوع‌های
$!A_1$، \dots، $!A_n$ از !!{formula} یکسان~$!A$ در $\delta$ است که در
آن
\begin{enumerate}
  \item $!A_1$ نتیجهٔ \Elim{\lor} یا \Elim{\lexists} نیست؛
  \item $!A_n$ مقدمهٔ فرعی \Elim{\lor} یا \Elim{\lexists} نیست؛
  \item برای $i = 1$، \dots،~$n-1$، $!A_i$ مقدمهٔ فرعی یک کاربرد از
  \Elim{\lor} یا \Elim{\lexists} است و $!A_{i+1}$ نتیجهٔ همان استنتاج
  است که پس از آن مقدمه می‌آید.
\end{enumerate}
\emph{طول} آن قطعه~$n$ است.
\end{defn}

هر قطعه‌ای هدف حذف (یعنی مسیر انحرافی) نیست. قطعه‌هایی که چنین‌اند
\emph{برش} نام دارند.

\begin{defn}
یک \emph{برش} قطعه‌ای است که در آن $!A_n$ مقدمهٔ اصلی یک استنتاج
!!a{elimination} است و یا $n=1$ و $!A_1 = !A_n$ نتیجهٔ یک استنتاج
!!a{introduction} یا نتیجهٔ~\FalseInt است، یا~$n>1$.
\end{defn}

توجه کنید که لازم نیست قطعهٔ برش با نتیجهٔ یک قاعدهٔ !!a{introduction}
آغاز شود. تنها شرط این است که به مقدمهٔ اصلی یک قاعدهٔ
!!a{elimination} ختم شود. اما برای آن‌که قطعه‌ای با طول~$1$ برش به شمار
آید، باید با نتیجهٔ یک قاعدهٔ !!{introduction} یا نتیجهٔ \FalseInt آغاز
شود.

\begin{defn}
\emph{رتبهٔ} یک برش~$\depth{!A}$ است. \emph{رتبهٔ
برش}~$\cutrank{\delta}$ در !!a{proof}~$\delta$ بیشینهٔ رتبهٔ برش‌های
موجود در~$\delta$ است و اگر برشی وجود نداشته باشد، برابر صفر است. یک
برش در~$\delta$ \emph{بیشینه} است هرگاه رتبه‌اش
~$\cutrank{\delta}$ باشد، یعنی رتبهٔ بیشینه داشته باشد. \emph{طول
برش}~$\delta$ مجموع طول همهٔ برش‌های بیشینهٔ آن است و اگر برشی وجود
نداشته باشد، برابر صفر است.
\end{defn}

\end{document}
